\section{The local formulas}\label{app:formulas}
This appendix recalls the finite expressions used in the comparison.
Their construction and estimates are the input results of
\cref{sec:inputs}. In the formulas below, $k$ is fixed and
$B=A+\mathscr C$.

\subsection{Extraction and reconstruction}
The extraction maps of \cite[\S3.1]{NormalizedContraction} use the
ordinary quadratic part
\[
 \mathscr E^{[2,0]}(X;A,\mathscr C)=\int\dd x\,\dd x'\,
 \bar\psi(x)^T K_X(A,\mathscr C;x,x')\psi(x').
\]
For a small $X$, choose one lift $\widetilde X$ and define
\begin{equation}\label{eq:rg-local-extractions}
 a_0(\mathscr E,X)
       =\operatorname{Vol}(X)^{-1}\mathscr E_{00}(X;0,0),\qquad
 a_2(\mathscr E,X)
       =\operatorname{Vol}(X)^{-1}
             \int_{\widetilde X}\dd x\int_{\widetilde X}\dd x'\,
                                                    K_X(0,0;x,x').
\end{equation}
Both cluster labels are averaged with measure $\eta_k^3/48$.
The second expression is independent of the lift: simultaneous
translation by a period multiplies its two arguments by the
same antiperiodic sign, whose product is one.

For an $M$-cube $\square$ put
\begin{equation}\label{eq:rg-counterterm-maps}
 \varepsilon_{\mathrm v}(\mathscr E,\square)
          =-\sum_{\substack{X\supset\square\\X\text{ small}}}
                                                     a_0(\mathscr E,X),
 \qquad
 M(\mathscr E,\square)
          =-\sum_{\substack{X\supset\square\\X\text{ small}}}
                                                     a_2(\mathscr E,X).
\end{equation}
Translations by $M\Z^3$ identify the small-polymer sets over
different cubes. They leave the two extracted quantities
unchanged. Thus the two sums do not depend on $\square$.
Write their values as $\varepsilon_{\mathrm v}(\mathscr E)$ and
$M(\mathscr E)$.

The imported scalarity assertion is $M(\mathscr E)=m(\mathscr E)1$.
On a small polymer the normalized remainder is
\cite[(17)]{NormalizedContraction}
\begin{equation}\label{eq:rg-normalized-remainder}
 (R\mathscr E)(X;A,\mathscr C)=\mathscr E(X;A,\mathscr C)
 -a_0(\mathscr E,X)\operatorname{Vol}(X)
 -\int_X\dd x\,\bar\psi(x)^Ta_2(\mathscr E,X)\psi(x).
\end{equation}
On other polymers it leaves the activity unchanged.
The operation $\Theta_{X,B}$ is the
shortest-chain quadratic adjustment of \cite[\S3.2]{NormalizedContraction}.
It averages over the shortest chains of $M$-cubes in $X$, using the
same six-path transports and pair coefficients $d_X^\alpha$.
We use the resulting $(R\mathscr E)^{\mathrm{nat}}$ from that definition.

Let $\Delta_y$ be the unit block centered at $y$.
For $c\in\Delta_y$, let $c_j(c)$ be the center of its block of side
$L^j\eta_k$. The averaging phase and its normalized average are
\[
 \varphi_A(y,c)=\exp\left(ie_k\eta_k\sum_{j=0}^{k-1}
          (\tau A)(c_{j+1}(c),c_j(c))\right),\qquad
 (\mathcal Q_k(A)f)(y)=\frac{\eta_k^3}{48}
       \sum_{c\in\Delta_y,g}\varphi_A(y,c)f(c,g).
\]
The reverse is the weighted transpose $\mathcal Q_k(-A)^T$.
For $y\in Y$ the two reconstruction rows of
\cite[Theorem~B]{NormalizedContraction} are as follows, where
$c_3=\pi x^{(3)}$.
\begin{equation}\label{eq:rg-reconstruction-two-pieces}
 \begin{split}
 (T_Y^{(1)}(A)\psi)(y)
 &=\int_{\Delta_y}\dd x\,\varphi_A(y,c)
    \Bigl\{(1+\mu/b_k)\psi(x)\\
 &\qquad-(b_k\eta_k)^{-1}\!\!\sum_{i:\,x^{(i)}\text{ over }Y}
           M_i(x)[U_A(x,x^{(i)})\psi(x^{(i)})-\psi(x)]\Bigr\},\\
 (T_Y^{(2)}(A)\chi)(y)
 &=-b_k^{-1}\eta_k^{\alpha-1}\int_{\Delta_y}\dd x\,
                1[c_3\notin Y]\varphi_A(y,c)M_3(x)\chi(x,x^{(3)}).
 \end{split}
\end{equation}
Rows outside $Y$ are zero. The first row contains ordinary
legs only over $Y$. The second row contains precisely the
type-$3$ pairs leaving $Y$, which are admissible boundary
pairs of length $\eta_k$. Their coefficients read $A$ only
on bonds with an endpoint in $\Delta_y$.
The barred rows are $\bar T_Y^{(i)}(B)=\mathsf C^{-1}T_Y^{(i)}(-B)\mathsf C$.
The substitution $\operatorname{sub}_{Y,B}$ replaces independent
unit generators $\Psi,\bar\Psi$ by these ordinary and pair rows and
fixes the old generators. This specifies the homomorphism used below.

\subsection{The determinant primitive and the first vertex}
Use the unprimed polymer resolvents $S_{k,y}(X,B)$ of
\cite[\S4.3]{FermionBlockManuscript}. Write
$\mathcal B_{k,y}=B_{k,y}$ for the finite block resolvent compared in
\cref{subsec:rgcov-finite-rows}, and $b_L=L^{-1}$.
At zero $F_0=D_B+\mu+b_LJP_LJ^*$ and $\Gamma_0=F_0^{-1}$.
The determinant primitive is \cite[(76)]{FermionBlockManuscript}:
\begin{align}
 \widetilde E_k(X,A)
    &=-ib_k^2\int_0^\infty\dd y\,
       \operatorname{Tr}\bigl[\gamma_5\mathcal B_{k,y}\mathcal Q_k
           S_{k,y}(X,A)\mathcal Q_k^\sharp\mathcal B_{k,y}\bigr]
                   \quad(k\ge1),\label{eq:local-determinant-definition}\\
 \widetilde E_0(X,A)
    &=\int_1^\infty\frac{\dd y}{y}
                     \operatorname{Tr}[F_0S_{0,y}(X,A)]
       -i\int_0^1\dd y\,\operatorname{Tr}[\gamma_5S_{0,y}(X,A)],\notag\\
 E_k^d(X,A)&=\tfrac14\bigl(\widetilde E_k(X,A)
       +\widetilde E_k(X,-A)
       +\overline{\widetilde E_k(X,\overline A)}
       +\overline{\widetilde E_k(X,-\overline A)}\bigr),\notag\\
 E_k^{\mathrm{det}}(X,A)&=E_k^d(X,A)-E_k^d(X,0).\notag
\end{align}
The traces use unit counting measure for coarse spinors and the
chamber measure $1/48$ at level zero. Absolute convergence and the
reblocked gain are part of the determinant input.

At zero write $q=e_0$ for the running charge. The local vertex of
\cite[(47)]{NonlinearStep} is
\begin{equation}\label{eq:rg-step-first-vertex}
 \begin{split}
 \mathcal V_0(Z,\square)=\frac12\int_\square\dd x\,
   \Big\{&
    \bar\psi(x)^TM_3(x)(e^{iqZ(c,c_3)}-1)
                                  [\psi(x)+\chi(x,x^{(3)})]\\
    &+[\bar\psi(x)+\bar\chi(x,x^{(3)})]^T
               M_3(x^{(3)})(e^{-iqZ(c,c_3)}-1)\psi(x)\Big\},
 \end{split}
\end{equation}
where $x=(c,g)$, $c_3=c+ge_1$ and
$\int_\square\dd x=48^{-1}\sum_{c\in\square}\sum_g$.
The pair $(x,x^{(3)})$ has $d_X=1$ and is admissible
for $\square$ of the second kind. Thus the vertex
has ordinary legs over the cube and precisely the
boundary pair legs allowed in the activity algebra.
Its coefficients are entire functions of $Z$ on
the bonds meeting the cube.


\subsection{The two finite localizations}
The definitions in this subsection are those of
\cite[\S5]{NonlinearStep}. For $\zeta\in\mathrm{Box}_k$ put
\begin{equation}\label{eq:rgnl-weakened-photon}
 Z(s')=H_k^{\mathrm{ph}}(s')C C_k^{1/2}(s')\zeta.
\end{equation}
The radii, with $e=e_k$, are
\begin{equation}\label{eq:rgnl-contour-radii}
 r_1=\tfrac12e^{-3/4+9\varepsilon/2},\qquad
 r_2=e^{-3/4+7\varepsilon},\qquad r_4=e^{-1/4+3\varepsilon}.
\end{equation}
The imported contour bounds give $2\le r_2,r_4\le r_1$ and
$A+tZ\in\widetilde R_k$, $B+tZ\in\widetilde G_k$ for $|t|\le r_1$.
The arguments of $H_k^\#$ are independent pairs
$(\Phi,\bar\Phi)$; in $H_k^\#\Phi$ the single symbol denotes
this pair. Write
$\psi^\#=(\psi,\chi;\bar\psi,\bar\chi)$ for the ordinary,
pair, barred ordinary and barred pair generators, respectively.
Then
\[
 \begin{aligned}
 H_k^\#(s'',B)\Phi
   ={}&(H_k(s'',B)\Phi,\delta_{\alpha,B}H_k(s'',B)\Phi;\\
       &
       \bar H_k(s'',B)\bar\Phi,
                       \delta_{\alpha,-B}\bar H_k(s'',B)\bar\Phi),\\
 \bar H_k(s'',B)&=\mathsf C^{-1}H_k(s'',-B)\mathsf C.
 \end{aligned}
\]
Define
$J_k^\#(s'',B,Z)=H_k^\#(s'',B+Z)-H_k^\#(s'',B)$ and
\[
 \widetilde\Gamma_k(s'',B)(W,\bar W)
                         =(\Gamma_k(s'',B)W,\bar W),\qquad
 W_k^\#(s'',B)=H_k^\#(s'',B)
                         \big(\Gamma_k(s'',B)W,\bar W\big).
\]
The notation $J_k,\bar J_k$ denotes the components without
H\"older quotient; $\mathrm{Wc},\overline{\mathrm{Wc}}$ denotes
the corresponding components of $W_k^\#$.
The auxiliary generators $W$ have weight
$\mathsf w_W=1$ for $k\ge1$ and $\sqrt{48}$ for $k=0$.

In the following localization, cubes and polymers are sets
of photon sites. Containment refers to all their lattice sites,
and two cubes meet when they share a lattice site.
Let $P_0(X)$ be the parameter cubes contained in $X$.
For a set $P$ of parameter cubes disjoint from $P_0(X)$, let $Y'(X,P)$ be the union of
the $M$-cubes meeting $X$ or a cube of $P_0(X)\cup P$.
For a continuous analytic family $G(s,X)$ define
\begin{align}
 \mathcal T_PG(X)
  &=\int_{[0,1]^P}\dd s_P\,
       \partial_PG(s,X)
       \big|_{s=1\text{ on }P_0(X),\,
                    s=0\text{ outside }P_0(X)\cup P},
                          \label{eq:rgnl-interpolation}\\
 \operatorname{Loc}G(Y)
  &=\sum_{(X,P):Y'(X,P)=Y}\mathcal T_PG(X).
                          \label{eq:rgnl-localization}
\end{align}
For a function analytic near the closed disk $|t|\le r$, $r>1$,
write
\[
 \mathcal C_rf=\frac1{2\pi i}\int_{|t|=r}
                                      \frac{f(t)}{t(t-1)}\,\dd t .
\]
This is $f(1)-f(0)$ and satisfies
$\|\mathcal C_rf\|\le(r-1)^{-1}\sup_{|t|=r}\|f(t)\|$.
Let $(m,\mathscr E)\in\operatorname{Ball}_k^x(4)$ and set
$m'=m+m(\mathscr E)$, $\mathscr E^r=R\mathscr E$.
For $k\ge1$ define
\begin{align}
 G_1(s,X)&=\mathcal C_{r_1}
                \langle\bar\Psi,M_k(X,B+tZ(s'))\Psi\rangle_1,
                                      \label{eq:rgnl-G1}\\
 G_2(s,X)&=\mathcal C_{r_2}\,
       \mathscr E^r(X;A+tZ(s'),\mathscr C,
                    \psi^\#+tJ_k^\#(s'',B,Z(s'))\Psi),
                                      \label{eq:rgnl-G2}\\
 G_3(s,\square)&=\mathcal C_{r_2}(-m')
       \langle\bar\psi+t\bar J_k(s'',B,Z(s'))\bar\Psi,
                         \psi+tJ_k(s'',B,Z(s'))\Psi\rangle_\square .
                                      \label{eq:rgnl-G3}
\end{align}
The symbols $M_k(X,\cdot)$ denote the local Dirac primitives.
A function defined for one cube is zero on the other polymers.
Put
\[
 E_i^*(Y)=\operatorname{Loc}G_i(Y),\qquad
 E_i(Y)=\operatorname{sub}_{Y,B}E_i^*(Y),\qquad
 E_{\le3}=\sum_{i=1}^3E_i.
\]
The substitution was defined above.
At $k=0$ there is no $\Psi$ and no substitution:
$G_1(s,\square)=\mathcal V_0(Z(s'),\square)$,
$G_2=G_3=0$, and $E_{\le3}=\operatorname{Loc}G_1$.
All parameter kinds occur in this first localization.

For the second localization, which uses only $\pi_k$, let
$v=W_k^\#(s'',B)$ and define
\begin{align}
 G_{\le3}^{\,W}(s'',X)&=E_{\le3}(X;A,\mathscr C,\zeta,
                                                      \psi^\#+v),
                                      \label{eq:rgnl-GW}\\
 G_4(s'',X)&=\mathcal C_{r_4}
                  \mathscr E^r(X;A,\mathscr C,\psi^\#+tv),\\
 G_5(s'',\square)&=\mathcal C_{r_4}(-m')
                 \langle\bar\psi+t\overline{\mathrm{Wc}},
                                   \psi+t\mathrm{Wc}\rangle_\square,\\
 F(Y)&=\operatorname{Loc}(G_{\le3}^{\,W}+G_4+G_5)(Y).
                                      \label{eq:rgnl-finite-activity}
\end{align}
At $k=0$ the term $G_4$ is zero and $m'=m$.


\subsection{Integration, logarithm and reblocking}
For the even activity $F$ just defined, set
\begin{align*}
 K(Y)&=\sum_{\substack{\{X_i\}\text{ distinct}\\\bigcup_iX_i=Y}}
             \prod_i(\exp F(X_i)-1),\qquad K(\varnothing)=1,\\
 H(Y)&=\int\dd\mu_I(W)\,\dd\widehat\mu_k(\zeta)\,K(Y),\\
 E_k^\#(Y)&=\sum_{n\ge1}\frac1{n!}
       \sum_{\bigcup_iY_i=Y}u^T(Y_1,\ldots,Y_n)\prod_iH(Y_i).
\end{align*}
Here $\widehat\mu_k$ is the product of the normal laws restricted
and normalized on $[-p_{0,k},p_{0,k}]$. The measure $\mu_I(W)$ is
the normalized identity-covariance Grassmann measure: the covariance
kernel is the inverse mode weight. The coefficient $u^T$ is the
sum of $(-1)^{|E(G)|}$ over connected spanning subgraphs of the
incompatibility graph of the tuple. Two polymers are incompatible
when their union is connected; in particular repetitions are incompatible.
Absolute convergence in the weighted coefficient norm is supplied by
\cite[Proposition~6.1]{NonlinearStep}.

For an outgoing polymer $Y$, put
\[
 E_k^*(Y;A',\mathscr C')=\sum_{\bar X^L=LY}
              S_LE_k^\#(X;A'_L,\mathscr C'_L),\qquad
 A'_L(b)=L^{-1/2}A'(b/L).
\]
The bar denotes the union of $LM$-cubes containing the $M$-cubes of $X$.
The map $S_L$ multiplies ordinary and pair legs by $L^{-1}$ and
$L^{-1-\alpha}$, respectively. The same sum defines $\mathcal L$
on any input activity. The scaled-domain inclusion used in comparison is
\begin{equation}\label{eq:rg-scaled-domain-inclusion}
 (A',\mathscr C')\in\widetilde R_{k+1}\times\mathfrak X_{k+1}
 \quad\Longrightarrow\quad
 (A'_L,\mathscr C'_L)\in\widetilde R_k^-\times\mathfrak X_k;
\end{equation}
it is \cite[(29)]{NormalizedContraction} under the assumed numerical conditions.
